% ============================================================
\chapter{MCMC --- Metropolis-Hastings}
\label{ch:metropolis}
% ============================================================

In 1953, at Los Alamos National Laboratory, Nicholas Metropolis, Arianna and Marshall Rosenbluth, and Augusta and Edward Teller published a paper on simulating molecules using a particular Monte Carlo method: instead of sampling the Boltzmann distribution directly (impossible in high dimensions), they constructed a Markov chain whose stationary distribution \emph{is} the target distribution. Twenty years later, W.\,Keith Hastings generalized the algorithm by allowing asymmetric proposal kernels. It would take until the 1990s for Bayesian statisticians --- notably Alan Gelfand and Adrian Smith --- to realize that this method solved \emph{their} fundamental problem: sampling from complex posterior distributions. The Metropolis--Hastings algorithm thus became the keystone of computational Bayesian inference, making tractable models that had seemed beyond reach.

\begin{intuition}[title=\textbf{Central idea}]
When the posterior has no closed form, we construct a Markov chain
whose stationary distribution is the posterior. The Metropolis--Hastings
algorithm is the foundational MCMC method: it proposes transitions and
accepts or rejects them according to a ratio that guarantees convergence
to the target.
\end{intuition}

% ------------------------------------------------------------
\section{Markov chain background}
% ------------------------------------------------------------

\begin{definition}[Markov chain]
A sequence $(\theta^{(t)})_{t \geq 0}$ is a \textbf{Markov chain} if
$\PP(\theta^{(t+1)} \in A \mid \theta^{(0)}, \ldots, \theta^{(t)})
= \PP(\theta^{(t+1)} \in A \mid \theta^{(t)})$ for all measurable $A$.
\end{definition}

\begin{definition}[Reversibility and detailed balance]
A transition kernel $K(\theta, \theta')$ satisfies \textbf{detailed
balance} with respect to $\pi$ if:
\[
  \pi(\theta)\,K(\theta, \theta') = \pi(\theta')\,K(\theta', \theta)
  \quad \forall \theta, \theta'.
\]
This implies that $\pi$ is the stationary distribution of the chain.
\end{definition}

\begin{theorem}[Ergodic theorem]
If the chain is irreducible and aperiodic with stationary distribution
$\pi$, then for every integrable function $g$:
\[
  \frac{1}{T}\sum_{t=1}^T g(\theta^{(t)}) \xrightarrow{\text{a.s.}}
  \int g(\theta)\,\pi(\theta)\,d\theta
  \quad \text{as } T \to \infty.
\]
\end{theorem}

% ------------------------------------------------------------
\section{Metropolis--Hastings algorithm}
% ------------------------------------------------------------

\begin{algorithme}[title=\textbf{Metropolis--Hastings algorithm}]
\begin{enumerate}
  \item Initialize $\theta^{(0)}$.
  \item For $t = 0, 1, 2, \ldots, T-1$:
  \begin{enumerate}
    \item Propose $\theta^* \sim q(\theta^* \mid \theta^{(t)})$.
    \item Compute the acceptance ratio:
    \[
      \alpha(\theta^{(t)}, \theta^*) = \min\!\left(1,\;
      \frac{\pi(\theta^* \mid x)\,q(\theta^{(t)} \mid \theta^*)}
           {\pi(\theta^{(t)} \mid x)\,q(\theta^* \mid \theta^{(t)})}\right).
    \]
    \item Draw $u \sim \text{Unif}(0, 1)$.
    \item If $u \leq \alpha$: set $\theta^{(t+1)} = \theta^*$ (accept).
          Else: set $\theta^{(t+1)} = \theta^{(t)}$ (reject).
  \end{enumerate}
\end{enumerate}
\end{algorithme}

\begin{theorem}[Validity of Metropolis--Hastings]
The Metropolis--Hastings kernel satisfies detailed balance with
respect to $\pi(\theta \mid x)$. If the chain is irreducible and
aperiodic, it converges to $\pi(\theta \mid x)$.
\end{theorem}

\begin{proof}
The transition kernel is:
\[
  K(\theta, \theta') = q(\theta' \mid \theta)\,\alpha(\theta, \theta')
  \quad (\theta' \neq \theta).
\]
To verify detailed balance, assume without loss of generality that
$\pi(\theta \mid x)\,q(\theta' \mid \theta) \leq
\pi(\theta' \mid x)\,q(\theta \mid \theta')$. Then
$\alpha(\theta, \theta') = 1$ and:
\begin{align*}
  \pi(\theta \mid x)\,K(\theta, \theta')
  &= \pi(\theta \mid x)\,q(\theta' \mid \theta) \\
  &= \pi(\theta' \mid x)\,q(\theta \mid \theta')
  \cdot \frac{\pi(\theta \mid x)\,q(\theta' \mid \theta)}
             {\pi(\theta' \mid x)\,q(\theta \mid \theta')} \\
  &= \pi(\theta' \mid x)\,q(\theta \mid \theta')\,\alpha(\theta', \theta) \\
  &= \pi(\theta' \mid x)\,K(\theta', \theta).
\end{align*}
\end{proof}

\begin{remark}
Only the ratio $\pi(\theta^*)/\pi(\theta^{(t)})$ is needed, so the
normalizing constant $m(x)$ cancels. It suffices to know the posterior
up to a constant.
\end{remark}

% ------------------------------------------------------------
\section{Special cases}
% ------------------------------------------------------------

\subsection{Metropolis algorithm (symmetric)}

\begin{definition}[Metropolis]
When the proposal is symmetric,
$q(\theta^* \mid \theta) = q(\theta \mid \theta^*)$, the ratio
simplifies to:
\[
  \alpha = \min\!\left(1,\; \frac{\pi(\theta^* \mid x)}{\pi(\theta^{(t)} \mid x)}\right).
\]
\end{definition}

\begin{example}
\textbf{Gaussian random walk.}
$q(\theta^* \mid \theta^{(t)}) = \mathcal{N}(\theta^{(t)}, \sigma_q^2 I)$.
Symmetry is immediate. The step size $\sigma_q$ must be calibrated: an
acceptance rate of $\approx 23\%$ is optimal in high dimensions.
\end{example}

\subsection{Independence sampler}

\begin{definition}[Independent Metropolis--Hastings]
The proposal does not depend on the current state:
$q(\theta^* \mid \theta^{(t)}) = q(\theta^*)$. The ratio becomes:
\[
  \alpha = \min\!\left(1,\;
  \frac{\pi(\theta^* \mid x)\,q(\theta^{(t)})}
       {\pi(\theta^{(t)} \mid x)\,q(\theta^*)}\right).
\]
\end{definition}

\begin{warning}[title=\textbf{Domination condition}]
The independence sampler works well only if $q$ dominates
$\pi(\cdot \mid x)$ in the sense that the tails of $q$ are at least as
heavy as those of $\pi$.
\end{warning}

% ------------------------------------------------------------
\section{Calibration and adaptation}
% ------------------------------------------------------------

\begin{keyformulas}[title=\textbf{Optimal acceptance rate}]
\begin{center}
\renewcommand{\arraystretch}{1.3}
\begin{tabular}{ll}
\toprule
\textbf{Dimension} & \textbf{Optimal rate} \\
\midrule
$d = 1$ & $\approx 44\%$ \\
$d \geq 2$ (Gaussian target) & $\approx 23.4\%$ \\
\bottomrule
\end{tabular}
\end{center}
Rule of thumb: adjust $\sigma_q$ to achieve an acceptance rate between
$20\%$ and $50\%$.
\end{keyformulas}

\begin{definition}[Adaptive Metropolis (AM)]
The AM algorithm adjusts the proposal covariance during burn-in:
\[
  \Sigma_q^{(t)} = \frac{2.38^2}{d}\,\hat{\Sigma}_t + \varepsilon I_d,
\]
where $\hat{\Sigma}_t$ is the empirical covariance of
$\theta^{(1)}, \ldots, \theta^{(t)}$.
\end{definition}

% ------------------------------------------------------------
\section{Convergence diagnostics}
% ------------------------------------------------------------

\begin{definition}[Trace plot]
A \textbf{trace plot} displays $\theta^{(t)}$ versus $t$. One looks
for ``white noise'' behavior (good mixing) with no visible trend or
correlation.
\end{definition}

\begin{definition}[$\hat{R}$ (Gelman--Rubin)]
The $\hat{R}$ diagnostic compares within-chain variance $W$ and
between-chain variance $B$ for $M$ chains:
\[
  \hat{R} = \sqrt{\frac{\hat{V}}{W}}, \qquad
  \hat{V} = \frac{n-1}{n}\,W + \frac{1}{n}\,B.
\]
Convergence if $\hat{R} < 1.01$ (strict criterion).
\end{definition}

\begin{definition}[Effective sample size (ESS)]
The \textbf{ESS} measures the number of equivalent independent samples:
\[
  \text{ESS} = \frac{T}{1 + 2\sum_{k=1}^\infty \rho_k},
\]
where $\rho_k$ is the lag-$k$ autocorrelation. The recommendation is
$\text{ESS} > 400$ per parameter.
\end{definition}

\begin{keyformulas}[title=\textbf{MCMC diagnostics checklist}]
\begin{enumerate}
  \item Check trace plots (no trends).
  \item $\hat{R} < 1.01$ for all parameters.
  \item $\text{ESS} > 400$ (or $> 100$ per chain).
  \item No divergences (for HMC/NUTS).
  \item Consistency across multiple chains.
\end{enumerate}
\end{keyformulas}

% ------------------------------------------------------------
\section{Python implementation}
% ------------------------------------------------------------

\begin{codeblock}[title=\textbf{Metropolis--Hastings from scratch}]
\begin{minted}{python}
import numpy as np
import matplotlib.pyplot as plt
from scipy import stats

# Target: Beta-Binomial posterior
n_data, s = 50, 15  # 15 successes out of 50
a_prior, b_prior = 1, 1

def log_posterior(theta):
    if theta <= 0 or theta >= 1:
        return -np.inf
    return (a_prior + s - 1) * np.log(theta) \
         + (b_prior + n_data - s - 1) * np.log(1 - theta)

# MH random walk
T = 50000
sigma_q = 0.1
samples = np.zeros(T)
samples[0] = 0.5
accepted = 0

for t in range(1, T):
    # Proposal
    theta_star = samples[t-1] + sigma_q * np.random.randn()
    # Ratio (Metropolis since symmetric)
    log_alpha = log_posterior(theta_star) \
              - log_posterior(samples[t-1])
    if np.log(np.random.rand()) < log_alpha:
        samples[t] = theta_star
        accepted += 1
    else:
        samples[t] = samples[t-1]

burn_in = 5000
samples_post = samples[burn_in:]
print(f"Acceptance rate: {accepted/T:.2%}")
print(f"Posterior mean: {samples_post.mean():.4f}")
print(f"Analytic: {(a_prior+s)/(a_prior+b_prior+n_data):.4f}")
\end{minted}
\end{codeblock}

\begin{codeblock}[title=\textbf{MCMC diagnostics}]
\begin{minted}{python}
fig, axes = plt.subplots(2, 2, figsize=(12, 8))

# Trace plot
axes[0, 0].plot(samples[:2000], alpha=0.7)
axes[0, 0].set_title("Trace plot (first 2000 iterations)")
axes[0, 0].set_xlabel("Iteration")
axes[0, 0].set_ylabel(r"$\theta$")

# Histogram vs analytic
theta_grid = np.linspace(0, 1, 200)
a_post, b_post = a_prior + s, b_prior + n_data - s
axes[0, 1].hist(samples_post, bins=50, density=True,
                alpha=0.5, label="MCMC")
axes[0, 1].plot(theta_grid,
                stats.beta.pdf(theta_grid, a_post, b_post),
                label="Analytic", color="red")
axes[0, 1].legend()
axes[0, 1].set_title("Posterior")

# Autocorrelation
from statsmodels.graphics.tsaplots import plot_acf
plot_acf(samples_post, lags=50, ax=axes[1, 0])
axes[1, 0].set_title("Autocorrelation")

# Running mean
running_mean = np.cumsum(samples_post) / \
               np.arange(1, len(samples_post) + 1)
axes[1, 1].plot(running_mean)
axes[1, 1].axhline((a_prior+s)/(a_prior+b_prior+n_data),
                    color="red", linestyle="--")
axes[1, 1].set_title("Running mean")
axes[1, 1].set_xlabel("Iteration")

plt.tight_layout()
plt.savefig("mcmc_diagnostics.pdf")
\end{minted}
\end{codeblock}

\begin{codeblock}[title=\textbf{Multidimensional MH with PyMC}]
\begin{minted}{python}
import pymc as pm
import arviz as az

# Bayesian regression model
np.random.seed(42)
n = 100
x = np.random.randn(n)
y = 2.5 + 1.3 * x + np.random.randn(n) * 0.8

with pm.Model() as model_reg:
    beta0 = pm.Normal("beta0", 0, 10)
    beta1 = pm.Normal("beta1", 0, 10)
    sigma = pm.HalfNormal("sigma", 5)
    mu = beta0 + beta1 * x
    y_obs = pm.Normal("y_obs", mu=mu, sigma=sigma,
                      observed=y)
    # Use Metropolis explicitly
    step = pm.Metropolis()
    trace_mh = pm.sample(20000, step=step,
                         tune=5000, random_seed=42)

# Diagnostics
az.plot_trace(trace_mh, var_names=["beta0", "beta1", "sigma"])
plt.savefig("trace_regression.pdf")

print(az.summary(trace_mh,
                 var_names=["beta0", "beta1", "sigma"]))
\end{minted}
\end{codeblock}

% ------------------------------------------------------------
\section{Exercises}
% ------------------------------------------------------------

\begin{exercise}
Implement the Metropolis random walk algorithm to sample from a
$t_3(0, 1)$ distribution. Study the effect of the step size $\sigma_q$
on the acceptance rate and ESS.
\end{exercise}

\begin{exercise}
Prove that the Metropolis--Hastings algorithm satisfies detailed
balance (complete proof).
\end{exercise}

\begin{exercise}
Implement the Adaptive Metropolis (AM) algorithm and compare with
standard Metropolis on a two-dimensional Gaussian posterior with strong
correlation ($\rho = 0.95$).
\end{exercise}

\begin{exercise}
For a two-component Gaussian mixture model, implement MH and observe
the multimodality problem. Propose solutions (parallel tempering,
large-jump proposals).
\end{exercise}

\begin{exercise}
\textbf{(Numerical)} Compare the efficiency (ESS per second) of
Metropolis--Hastings with NUTS (PyMC default) on a logistic regression
model with $p = 10$ predictors.
\end{exercise}
